% Capa
\imprimircapa

% Folha de rosto
\imprimirfolhaderosto

% Resumo
\setlength{\absparsep}{18pt} % ajusta o espaçamento dos parágrafos do resumo
\begin{resumo}
	\SingleSpacing
	Este trabalho apresenta um estudo de caso sobre a estimativa da constante matemática \texorpdfstring{$\pi$}{pi} por meio do método estocástico de Monte Carlo. O objetivo principal é analisar a convergência do algoritmo e avaliar o erro numérico (absoluto, relativo e relativo percentual) em função do aumento do número de amostras. A simulação foi implementada na linguagem Python, consistindo no sorteio aleatório de coordenadas bidimensionais e na verificação de sua incidência geométrica em um quarto de circunferência unitária inscrito em um quadrado. Os resultados obtidos para uma amostragem de 10.000 iterações demonstraram a eficácia do método, alcançando um erro relativo percentual inferior a 0,5\%. Conclui-se que o método de Monte Carlo, fundamentado na Lei dos Grandes Números, é uma ferramenta robusta para aproximações numéricas, embora dependa da qualidade do gerador de números pseudoaleatórios e de um volume amostral expressivo para atenuar as flutuações estocásticas.
	
	\vspace{\onelineskip}
	\noindent
	\textbf{Palavras-chave}: Método de Monte Carlo. Constante Pi. Processos Estocásticos. Análise de Erro. Simulação Computacional.
\end{resumo}

% Abstract
\begin{resumo}[Abstract]
	\SingleSpacing
	This paper presents a case study on estimating the mathematical constant \texorpdfstring{$\pi$}{pi} using the stochastic Monte Carlo method. The main objective is to analyze the algorithm's convergence and evaluate the numerical error (absolute, relative, and percentage relative error) as a function of the increasing number of samples. The simulation was implemented in Python, consisting of the random drawing of two-dimensional coordinates and the verification of their geometric incidence in a unit quarter-circle inscribed in a square. The results obtained for a sampling of 10,000 iterations demonstrated the effectiveness of the method, achieving a relative percentage error of less than 0.5\%. It is concluded that the Monte Carlo method, based on the Law of Large Numbers, is a robust tool for numerical approximations, although it depends on the quality of the pseudorandom number generator and a significant sample volume to mitigate stochastic fluctuations.
		
	\vspace{\onelineskip}
	\noindent
	\textbf{Keywords}: Monte Carlo Method. Pi Constant. Stochastic Processes. Error Analysis. Computational Simulation.
\end{resumo}

{%hidelinks
	\hypersetup{hidelinks}
	
	% inserir lista de figuras
	\pdfbookmark[0]{\listfigurename}{lof}
	\listoffigures*
	\cleardoublepage
	
	% inserir lista de tabelas
	\pdfbookmark[0]{\listtablename}{lot}
	\listoftables*
	\cleardoublepage
	
	% inserir lista de abreviaturas e siglas (devem ser declarados no preambulo)
	\ifnotempty{\verificasiglas}{
	\imprimirlistadesiglas
	}
	
	% inserir lista de símbolos (devem ser declarados no preambulo)
	\ifnotempty{\verificasimbolos}{
	\imprimirlistadesimbolos
	}
	
	% inserir o sumario
	\pdfbookmark[0]{\contentsname}{toc}
	\tableofcontents*
	\cleardoublepage
	
}%hidelinks